\section*{ЗАКЛЮЧЕНИЕ}
\addcontentsline{toc}{section}{ЗАКЛЮЧЕНИЕ}

В результате выполнения выпускной квалификационной работы была достигнута поставленная цель --- разработана программно-информационная система регистрации данных пользователей, предназначенная для автоматизации основных бизнес-процессов организации сервисного обслуживания автомобилей. Разработанное решение реализовано в виде современного web-приложения и обеспечивает централизованное хранение данных, разграничение прав доступа, контроль жизненного цикла заявок и ведение журнала изменений.

В ходе выполнения работы был проведён анализ предметной области, позволивший выявить существующие проблемы организации процессов регистрации и обработки данных пользователей. Было установлено, что использование электронных таблиц, бумажных носителей и разрозненных программных решений приводит к дублированию информации, увеличению времени обработки обращений, затрудняет контроль действий сотрудников и повышает вероятность возникновения ошибок, связанных с человеческим фактором.

Проведённый анализ существующих программных решений показал, что большинство представленных на рынке систем либо ориентированы на крупные организации и требуют значительных затрат на внедрение и сопровождение, либо содержат избыточную функциональность, не соответствующую потребностям небольших и средних предприятий сервисного обслуживания. Кроме того, многие решения имеют ограниченные возможности адаптации под особенности конкретной организации. Это подтвердило актуальность разработки собственной программно-информационной системы, учитывающей специфику рассматриваемой предметной области.

В процессе выполнения работы были решены все поставленные задачи:

\begin{enumerate}
	\item проведён анализ предметной области и выявлены бизнес-процессы, подлежащие автоматизации;
	
	\item определены функциональные и нефункциональные требования к разрабатываемой программной системе;
	
	\item спроектирована архитектура программного решения на основе доменного подхода с разделением ответственности между уровнями пользовательского интерфейса, бизнес-логики и доступа к данным;
	
	\item разработана структура системы, включающая модули аутентификации и авторизации, управления клиентами, автомобилями, заявками на обслуживание, пользователями системы и журналом изменений;
	
	\item реализованы механизмы взаимодействия с облачной платформой Supabase для хранения и обработки данных;
	
	\item реализована ролевая модель доступа, обеспечивающая разграничение полномочий пользователей в зависимости от их должностных обязанностей;
	
	\item разработан алгоритм управления жизненным циклом заявок, основанный на модели конечного автомата и исключающий возникновение недопустимых переходов между состояниями;
	
	\item реализованы средства аудита действий пользователей, обеспечивающие возможность контроля и анализа выполняемых операций;
	
	\item проведено функциональное и системное тестирование разработанной системы, подтвердившее корректность её работы и соответствие предъявляемым требованиям.
\end{enumerate}

Разработанная программно-информационная система построена с использованием современных web-технологий: фреймворка Vue~3, языка программирования TypeScript, библиотеки управления состоянием Pinia и облачной платформы Supabase. Применение данных технологий позволило создать масштабируемое, удобное в сопровождении и расширении программное решение, не требующее установки специализированного программного обеспечения на рабочих местах пользователей.

Использование архитектурного подхода, основанного на разделении системы на независимые домены, обеспечило низкую связанность между компонентами приложения и повысило удобство дальнейшего развития проекта. Каждая функциональная область системы была выделена в отдельный модуль, содержащий пользовательские интерфейсы, бизнес-логику, механизмы доступа к данным и описание используемых типов. Подобная организация исходного кода способствует повышению читаемости проекта, упрощает его сопровождение и снижает трудозатраты при внедрении новых возможностей.

Особое внимание в работе было уделено вопросам обеспечения информационной безопасности. Реализованная многоуровневая система контроля доступа включает разграничение прав на уровне пользовательского интерфейса, проверку маршрутов приложения, контроль бизнес-правил в программном коде и политики Row Level Security на уровне базы данных. Такой подход позволяет обеспечить защиту данных даже при попытках обхода механизмов клиентской части приложения.

Проведённое системное тестирование показало устойчивую работу программной системы при выполнении основных пользовательских сценариев. Были успешно проверены процессы авторизации пользователей, регистрации и редактирования клиентов, учёта автомобилей, создания и обработки заявок, управления учётными записями сотрудников, а также функционирование журнала изменений. Результаты тестирования подтвердили корректность реализованных алгоритмов и соответствие фактического поведения системы ожидаемым результатам.

Практическая значимость выполненной работы заключается в возможности использования разработанной системы в деятельности организаций сервисного обслуживания автомобилей для автоматизации процессов регистрации и обработки данных пользователей. Внедрение системы позволит сократить время выполнения рутинных операций, уменьшить вероятность возникновения ошибок, повысить прозрачность деятельности сотрудников и обеспечить централизованный доступ к актуальной информации.

Разработанное программное решение обладает потенциалом дальнейшего развития. В качестве перспективных направлений совершенствования системы можно выделить реализацию системы уведомлений клиентов посредством электронной почты и SMS-сообщений, интеграцию с внешними сервисами поставщиков запасных частей, формирование расширенной аналитической отчётности, внедрение механизмов резервного копирования и восстановления данных, а также разработку мобильной версии приложения для сотрудников организации.

Таким образом, поставленная цель выпускной квалификационной работы достигнута в полном объёме. Разработанная программно-информационная система регистрации данных пользователей соответствует сформулированным требованиям, успешно прошла тестирование и может быть рекомендована к практическому внедрению и дальнейшему развитию в соответствии с потребностями организаций.